\section{Discussion}
\label{sec:discussion}

\subsection{A mechanism can be real and still not pay}

The result that organises this paper is not that a frequency-aware objective
failed. It is that the objective \emph{succeeded} at the thing it was designed
to do and lost anyway. It closes \correction~\si{\decibel} of a
\deficitHigh~\si{\decibel} high-frequency deficit; the correction is
attributable to the component its design credits, since removing the temporal
schedule removes the whole effect (\S\ref{sec:mechanism}); and the correction
is worth $\boostValue$~FID when applied in isolation (\S\ref{sec:worth}). The
objective nonetheless lands \shortfall~FID short of what that predicts.

This is a different failure from the one usually reported. A method that does
not work is uninteresting; a method that does exactly what it claims and is
still outperformed by its baseline says something about the relationship
between the mechanism and the goal. Here the relationship is that the mechanism
is not free: reallocating the error budget towards high frequencies means
allocating less of it everywhere else, and by Eq.~\ref{eq:parseval} that is
precisely what the reallocation must do. What the objective gives up is not
visible in the spectrum, because the spectrum is the one thing it was
optimising.

\subsection{Post-processing as a baseline}

That the same spectral correction can be obtained by a Fourier operation on
finished samples --- outperforming the trained objective --- is not offered as
a method. It is offered as a control that this literature does not use and, we
think, should.

Any objective whose stated mechanism is a change to an image statistic admits
the same test: apply that change to the baseline's outputs directly and see
what it is worth. If the trained method does not clear the post-processed
baseline, the training-time machinery is not paying for itself, whatever the
ablations show. The test is cheap, it requires no additional training, and it
would have flagged the objective audited here before any of the sweeps in
\S\ref{sec:results} were run.

\subsection{What the metrics can and cannot register}

Our two measurement hazards share a structure: in both, a quantity was
estimated correctly under one set of conditions and read as though it held
under another.

The KID reversal (\S\ref{sec:replication}) is the familiar version --- a
bespoke estimator, correct in its own terms, disagreeing in sign with the
standard one. It is worth reporting because the disagreement was the study's
only positive result, and because the failure is invisible unless a second
implementation is run.

The calibration error (\S\ref{sec:worth}) is subtler and, we suspect, more
common. Measuring how much FID responds to a perturbation of real images
measures the response \emph{at the minimum}, where the first-order term
vanishes by construction and the metric is quadratically insensitive. It is a
natural experiment to design and it gives an answer roughly thirty times too
small. Sensitivity has to be probed where the model actually sits. We report
this because we drew the wrong conclusion from it first, and because the same
design suggests itself to anyone asking the same question.

Neither hazard is an argument against these metrics. Both are arguments for
asking what a result \emph{should} have been before deciding what it means.

\subsection{Implications for frequency-aware training}

We audit one objective, and we are careful not to generalise the empirical
result past it. But three parts of the analysis are not specific to it.

Equation~\ref{eq:parseval} constrains the whole family: on an orthonormal
basis, sub-band weighting redistributes a fixed budget, so every gain is
someone else's loss and the question is always whether the trade is favourable.
The confound of \S\ref{sec:confound} arises whenever a weighting depends on
$\sigma_t$ without its total being held fixed, which is the common case; the
protocol that separates the axes costs nothing to adopt. And the
post-processing control of \S\ref{sec:worth} applies to any method that names
the statistic it intends to change.

Whether other members of the family clear these bars is an empirical question
we have not answered. We would expect the answer to depend on resolution: at
$32\times32$ the high-frequency band carries little of the signal that
Inception features encode, and the trade a frequency-aware objective makes may
look different where it carries more.
